%% kapitel/evaluation.tex
\section{Evaluation}
\label{sec:evaluation}

\section{Testumgebung und Datensätze}
Die Evaluation der \textit{Hackintosh Companion App} erfolgte primär auf einem physischen Testgerät des Typs \textbf{Samsung Galaxy A3 (2016)} unter Android 13. Zur Überprüfung der Skalierbarkeit und Performance wurden 50 Test-Datensätze generiert, die komplexe Hardware-Konfigurationen (inkl. Listen von Kexts und ACPI-Patches) abbilden.

\section{Performance-Analyse}

\subsection{Datenbank-Performance (SQLite)}
Die lokale Datenhaltung mittels SQLite zeigt eine hohe Effizienz. Lesezugriffe auf die Gerätedatenbank sowie die Filterung nach spezifischen Hardware-IDs erfolgen in weniger als 15ms. Die Implementierung von Indizes auf den Feldern \texttt{name} und \texttt{device\_id} gewährleistet auch bei steigender Datenmenge eine flüssige Darstellung in der \texttt{ListView}.

\subsection{Grafik- und Animations-Performance}
Ein kritischer Punkt war die Ausführung der 3D-Visualisierung und der Boot-Animation auf der limitierten Hardware des Samsung A3.
\begin{itemize}
    \item \textbf{Boot-Animation:} Durch die Nutzung asynchroner Timer (\texttt{Future.delayed}) wird der Main-Thread entlastet, was zu einer ruckelfreien Darstellung der BIOS-Logs führt.
    \item \textbf{3D-Rendering:} Die Deaktivierung des ressourcenintensiven \textit{Impeller}-Renderers zugunsten von \textit{Skia} resultierte in einer stabilen Framerate von ca. 50--60 FPS während der Interaktion mit dem Laptop-Modell.
\end{itemize}

\section{Funktionale Ergebnisse}
Die Evaluation der Kernfunktionen lieferte folgende Ergebnisse:
\begin{itemize}
    \item \textbf{Offline-Verfügbarkeit:} Der \textit{Offline-First-Ansatz} wurde erfolgreich verifiziert. Alle Gerätedaten sind ohne aktive Internetverbindung sofort nach dem App-Start verfügbar.
    \item \textbf{QR-Code-Erkennung:} Der \texttt{mobile\_scanner} erkennt komplexe JSON-Strings innerhalb von 1,2 Sekunden unter normalen Lichtbedingungen.
    \item \textbf{Cloud-Synchronisation:} Der Datenaustausch via REST-API funktioniert stabil. Die Fehlerbehandlung (\textit{Error Handling}) fängt Verbindungsabbrüche ohne Absturz der Anwendung ab.
\end{itemize}

\section{Zusammenfassende Beurteilung}
Die App erfüllt alle im Konzept definierten Anforderungen. Besonders hervorzuheben ist die Kombination aus technischer Informationstiefe und intuitiver Benutzerführung. Die Entscheidung für Flutter als Framework ermöglichte es, eine visuell ansprechende Applikation zu entwickeln, die trotz der hardwarenahen Features (3D, Kamera) auch auf älteren Endgeräten eine hohe User Experience bietet.